
\section{\BookRecordIndex{T0283}求值映射的\BookRecordIndex{T0166}横截性}\label{sec:6-3}
\zhindex{qiuzhiyingshedehengjiexing}{求值映射的横截性}

本节证明命题~\ref{prop:6-2-8}：对任意树 $T$，\BookRecordIndex{T0283}求值映射
\[
  \operatorname{ev}^E:
  \mathcal M^*(\{A_\alpha\};\mathcal J^\ell)\times Z(T)
  \longrightarrow M^E
\]
与对角线 $\Delta^E$ 横截。事实证明，在比树更一般的对象上证明这个结果会很方便；也就是说，
不再假设任意两个顶点都能由一条\BookRecordIndex{T0064}边链连接。若有限集 $T$ 配备的关系
$E\subset T\times T$ 满足附录 D.2 的（对称性）、（反自反性）与（无圈）条件，则称它为
一个\emph{\BookRecordIndex{T0297}森林}；这里不要求（连通性）公理。若没有\BookRecordIndex{T0064}边链连接 $\alpha$ 与 $\beta$，记
$[\alpha,\beta]:=\varnothing$。本章所有定义都以显然的方式推广到\BookRecordIndex{T0297}森林。例如，与
式~\eqref{eq:6-2-4} 相同，\BookRecordIndex{T0184}简单稳定映射的\BookRecordIndex{T0268}普适模空间是\BookRecordIndex{T0269}普适求值映射下对角线的原像
\[
  \widetilde{\mathcal M}^*_{0,T}(\{A_\alpha\};\mathcal J^\ell)
  =\bigl(\operatorname{ev}^E\bigr)^{-1}(\Delta^E)
  \subset
  \mathcal M^*(\{A_\alpha\};\mathcal J^\ell)\times Z(T),
\]
而\BookRecordIndex{T0413}重参数化群 $G_T$ 的定义与先前相同。只要 $\operatorname{ev}^E$ 与带标号\BookRecordIndex{T0297}森林
$(T,E,\Lambda)$ 的\BookRecordIndex{T0061}边对角线 $\Delta^E$ 横截，空间
$\widetilde{\mathcal M}^*_{0,T}(\{A_\alpha\};\mathcal J^\ell)$ 就是
$C^{\ell-1}$ \BookRecordIndex{T0053}巴拿赫流形。在这种情形下，$G_T$ 通过 $C^{\ell-1}$ 微分同胚自由且恰当地
作用在这个空间上，因为对每个
$(\mathbf u,\mathbf z,J)\in
\widetilde{\mathcal M}^*_{0,T}(\{A_\alpha\};\mathcal J^\ell)$，每个 $u_\alpha$ 都属于
$W^{\ell+1,p}$ 类。因此，只要 $\operatorname{ev}^E$ 与 $\Delta^E$ 横截，商
\[
  \mathcal M^*_{0,T}(\{A_\alpha\};\mathcal J^\ell)
  :=\widetilde{\mathcal M}^*_{0,T}
  (\{A_\alpha\};\mathcal J^\ell)/G_T
\]
就是 $C^{\ell-1}$ \BookRecordIndex{T0053}巴拿赫流形。

命题~\ref{prop:6-2-8} 的证明建立在命题~\ref{prop:3-4-2} 的如下对应版本之上。与
第~\ref{sec:3-4} 节一样，一个较弱的陈述有初等证明：对角线的补集完全由\BookRecordIndex{T0283}求值映射
$\operatorname{ev}_{ij}$ 的\BookRecordIndex{T0409}正则值组成。该证明利用
$\operatorname{Diff}(M,\omega)$ 在\BookRecordIndex{T0268}普适模空间
$\mathcal M^*_{0,T}(\{A_\alpha\};\mathcal J^\ell)$ 上的作用；细节留给读者。这里说明一个复杂
得多的论证，它用来证明对角线上的每一点也都是\BookRecordIndex{T0409}正则值。

\begin{theorem}\label{thm:6-3-1}
设 $T=(T,E,\Lambda)$ 是一个 $k$-带标号\BookRecordIndex{T0297}森林，并设映射
\[
  T\longrightarrow H_2(M;\mathbb Z),
  \qquad \alpha\longmapsto A_\alpha
\]
满足稳定性条件~\eqref{eq:6-1-1}。假设\BookRecordIndex{T0065}边求值映射
\[
  \operatorname{ev}^E:
  \mathcal M^*(\{A_\alpha\};\mathcal J^\ell)\times Z(T)
  \longrightarrow M^E
\]
与 $\Delta^E$ 横截。设 $i,j\in\{1,\ldots,k\}$，且顶点
$\alpha_i,\alpha_j$ 属于 $T$ 的不同连通分支，即
$[\alpha_i,\alpha_j]=\varnothing$。则 $M\times M$ 中的每一点都是\BookRecordIndex{T0283}求值映射
\[
  \operatorname{ev}_{ij}:
  \mathcal M^*_{0,T}(\{A_\alpha\};\mathcal J^\ell)
  \longrightarrow M\times M
\]
的\BookRecordIndex{T0409}正则值；该映射定义为
\[
  \operatorname{ev}_{ij}([\mathbf u,\mathbf z,J])
  :=\bigl(u_{\alpha_i}(z_i),u_{\alpha_j}(z_j)\bigr).
\]
\end{theorem}

\begin{proof}
由命题~\ref{prop:6-2-7}，
$\mathcal M^*(\{A_\alpha\};\mathcal J^\ell)$ 是\BookRecordIndex{T0053}巴拿赫流形；由假设，\BookRecordIndex{T0269}普适求值映射
$\operatorname{ev}^E$ 与 $\Delta^E$ 横截。于是，\BookRecordIndex{T0380}隐函数定理说明\BookRecordIndex{T0184}简单稳定映射的\BookRecordIndex{T0268}普适模空间
\[
  \widetilde{\mathcal M}^*_{0,T}(\{A_\alpha\};\mathcal J^\ell)
  =\bigl(\operatorname{ev}^E\bigr)^{-1}(\Delta^E)
\]
\sourcefootnote{原书此处把\BookRecordIndex{T0246}模空间的同调类参数印作总类 $A$；但定理固定的是整族
$\{A_\alpha\}$，本节前后的\BookRecordIndex{T0268}普适模空间也都以这族分支类为参数，故改正。}
是\BookRecordIndex{T0053}巴拿赫流形。它在点 $(\mathbf u,\mathbf z,J)$ 处的切空间由所有元组
\[
  (\boldsymbol\xi,\boldsymbol\zeta,Y)
  =\bigl(\{\xi_\alpha\}_{\alpha\in T},
    \{\zeta_{\alpha\beta}\}_{\alpha E\beta},
    \{\zeta_i\}_{1\le i\le k},Y\bigr)
\]
组成，其中
$\xi_\alpha\in W^{\ell,p}(S^2,u_\alpha^*TM)$、
$\zeta_{\alpha\beta}\in T_{z_{\alpha\beta}}S^2$、
$\zeta_i\in T_{z_i}S^2$、$Y\in T_J\mathcal J^\ell$，并且满足
\begin{equation}\label{eq:6-3-1}
  D_{u_\alpha}\xi_\alpha
  +\frac12Y(u_\alpha)\circ du_\alpha\circ j=0,
\end{equation}
以及对所有 $\alpha E\beta$，
\begin{equation}\label{eq:6-3-2}
  \xi_\alpha(z_{\alpha\beta})
  +du_\alpha(z_{\alpha\beta})\zeta_{\alpha\beta}
  =\xi_\beta(z_{\beta\alpha})
  +du_\beta(z_{\beta\alpha})\zeta_{\beta\alpha}.
\end{equation}
注意，式~\eqref{eq:6-3-2} 正是条件
$u_\alpha(z_{\alpha\beta})=u_\beta(z_{\beta\alpha})$ 的导数。\BookRecordIndex{T0283}求值映射
$\operatorname{ev}_{ij}$ 在 $(\mathbf u,\mathbf z,J)$ 处的微分为
\begin{equation}\label{eq:6-3-3}
\begin{aligned}
  &d\operatorname{ev}_{ij}(\mathbf u,\mathbf z,J)
  (\boldsymbol\xi,\boldsymbol\zeta,Y)\\
  &\qquad=
  \bigl(\xi_{\alpha_i}(z_i)+du_{\alpha_i}(z_i)\zeta_i,
        \xi_{\alpha_j}(z_j)+du_{\alpha_j}(z_j)\zeta_j\bigr).
\end{aligned}
\end{equation}
我们证明：对每个
$(\mathbf u,\mathbf z,J)\in
\widetilde{\mathcal M}^*_{0,T}(\{A_\alpha\};\mathcal J^\ell)$，这个微分都满射。由 $i,j$
之间的对称性，只需证明子空间 $T_{u_{\alpha_i}(z_i)}M\times\{0\}$ 包含在
$d\operatorname{ev}_{ij}(\mathbf u,\mathbf z,J)$ 的像中。为此，固定向量
$v_i\in T_{u_{\alpha_i}(z_i)}M$。我们要证明 $(v_i,0)$ 属于该微分的像。

先假设 $u_{\alpha_i}$ 非常值。由命题~\ref{prop:2-5-1} 和命题~\ref{prop:2-4-4}，
存在 $u_{\alpha_i}$ 的\BookRecordIndex{T0094}单射点 $w_i\in S^2$，使得对每个
$\beta\ne\alpha_i$，
$u_{\alpha_i}(w_i)\notin u_\beta(S^2)$。选取 $\varepsilon>0$，使得对每个
$\beta\in T$，
\[
  \beta\ne\alpha_i
  \quad\Longrightarrow\quad
  B_\varepsilon(u_{\alpha_i}(w_i))\cap u_\beta(S^2)=\varnothing.
\]
于是，由引理~\ref{lem:3-4-3}，存在一对 $(\xi_{\alpha_i},Y)$，它对
$\alpha=\alpha_i$ 满足式~\eqref{eq:6-3-1}，并且
\[
  \xi_{\alpha_i}(z_i)=v_i,
  \qquad
  \xi_{\alpha_i}(z_{\alpha_i\beta})=0
  \qquad
  \operatorname{supp}Y\subset B_\varepsilon(u_{\alpha_i}(w_i))
\]
\sourcefootnote{原书把这里的向量场误印作带两个顶点下标的
$\xi_{\alpha_i\beta}$；切向量元组只定义了按顶点标号的 $\xi_\alpha$，故改为
$\xi_{\alpha_i}$。}
对每个满足 $\alpha_iE\beta$ 的 $\beta\in T$ 成立。因此，对
$\beta\ne\alpha_i$，有 $Y\circ u_\beta\equiv0$。在元组
$(\boldsymbol\xi,\boldsymbol\zeta,Y)$ 中令
$\boldsymbol\zeta:=0$，并对 $\beta\ne\alpha_i$ 令 $\xi_\beta:=0$，则它满足
式~\eqref{eq:6-3-1} 与式~\eqref{eq:6-3-2}。此外，
$\alpha_j\ne\alpha_i$，故
\[
  d\operatorname{ev}_{ij}(\mathbf u,\mathbf z,J)
  (\boldsymbol\xi,\boldsymbol\zeta,Y)=(v_i,0).
\]

现在设 $u_{\alpha_i}$ 是常值映射，考虑\BookRecordIndex{T0417}子树
\[
  T_i:=\bigl\{\alpha\in T\bigm|
  [\alpha_i,\alpha]\ne\varnothing,
  \ A_\beta=0\text{ 对每个 }\beta\in[\alpha_i,\alpha]\bigr\}.
\]
它由所有能通过一条\BookRecordIndex{T0382}幽灵分支链与 $\alpha_i$ 相连的顶点组成。对每个 $\alpha\in T_i$，都有
$u_\alpha\equiv u_{\alpha_i}(z_i)$。记
\[
  V_i:=\{\beta\in T\setminus T_i\mid
  \text{存在 }\alpha\in T_i\text{ 使 }\alpha E\beta\}
\]
为与 $T_i$ 相邻的顶点集合。于是，对每个 $\beta\in V_i$ 都有 $A_\beta\ne0$；并且存在
唯一映射 $f:V_i\to T_i$，使得对每个 $\beta\in V_i$ 都有 $f(\beta)E\beta$。
此外，由命题~\ref{prop:2-5-1} 和命题~\ref{prop:2-4-4}，对每个
$\beta\in V_i$，都存在 $u_\beta$ 的\BookRecordIndex{T0094}单射点 $w_\beta\in S^2$，使得对每个
$\gamma\in T\setminus\{\beta\}$，
$u_\beta(w_\beta)\notin u_\gamma(S^2)$。选取 $\varepsilon>0$，使对所有
$\beta,\gamma\in T$ 都有
\[
  \beta\in V_i,\quad\beta\ne\gamma
  \quad\Longrightarrow\quad
  B_\varepsilon(u_\beta(w_\beta))\cap u_\gamma(S^2)=\varnothing.
\]

现在设 $\beta\in V_i$。由引理~\ref{lem:3-4-3}，存在一对
$(\xi_\beta,Y_\beta)$，它满足式~\eqref{eq:6-3-1}，并且
\[
  \xi_\beta(z_{\beta f(\beta)})=v_i,
  \qquad
  \xi_\beta(z_{\beta\gamma})=0,
  \qquad
  \operatorname{supp}Y_\beta
  \subset B_\varepsilon(u_\beta(w_\beta))
\]
对每个满足 $\beta E\gamma$ 的
$\gamma\in T\setminus\{f(\beta)\}$ 成立。因此，当 $\beta\in V_i$ 且
$\gamma\ne\beta$ 时，$Y_\beta\circ u_\gamma\equiv0$。在元组
$(\boldsymbol\xi,\boldsymbol\zeta,Y)$ 中，对 $\beta\in V_i$ 采用刚刚选定的
$\xi_\beta$，并定义
\[
  \boldsymbol\zeta:=0,
  \qquad
  \xi_\alpha:=
  \begin{cases}
    v_i,&\alpha\in T_i,\\
    0,&\alpha\in T\setminus(T_i\cup V_i),
  \end{cases}
  \qquad
  Y:=\sum_{\beta\in V_i}Y_\beta.
\]
于是，对每个 $\alpha\in T\setminus V_i$，都有
$\xi_\alpha\in\ker D_{u_\alpha}$ 且 $Y\circ u_\alpha=0$；对每个
$\beta\in V_i$，则有 $Y\circ u_\beta=Y_\beta\circ u_\beta$。因此，元组
$(\boldsymbol\xi,\boldsymbol\zeta,Y)$ 满足式~\eqref{eq:6-3-1} 与
式~\eqref{eq:6-3-2}。此外，$\alpha_j\notin T_i\cup V_i$，故
\[
  d\operatorname{ev}_{ij}(\mathbf u,\mathbf z,J)
  (\boldsymbol\xi,\boldsymbol\zeta,Y)=(v_i,0).
\]
这证明了定理~\ref{thm:6-3-1}。
\end{proof}

\begin{proof}[命题~\ref{prop:6-2-8} 的证明]
我们证明一个更强的结果：对每个\BookRecordIndex{T0297}森林 $(T,E,\Lambda)$，
$\operatorname{ev}^E$ 都与 $\Delta^E$ 横截。当\BookRecordIndex{T0297}森林没有边时，断言为空。设 $N$ 是非负整数，
并归纳假设该断言已对至多有 $N$ 条边的所有\BookRecordIndex{T0297}森林成立。设 $T$ 是有 $N+1$ 条边的
$k$-带标号\BookRecordIndex{T0297}森林。任选一条边 $\alpha E\beta$，切断这条边，并把两个点
$z_{\alpha\beta},z_{\beta\alpha}$ 加入\BookRecordIndex{T0066}标记点族，由此得到带标号\BookRecordIndex{T0297}森林
$T'=(T',E',\Lambda')$。它定义为
\[
  T':=T,
  \qquad
  E':=E\setminus\{(\alpha,\beta),(\beta,\alpha)\},
\]
以及
\[
  \Lambda'_\gamma:=
  \begin{cases}
    \Lambda_\gamma,&\gamma\notin\{\alpha,\beta\},\\
    \{\alpha\beta\}\cup\Lambda_\gamma,&\gamma=\alpha,\\
    \{\beta\alpha\}\cup\Lambda_\gamma,&\gamma=\beta.
  \end{cases}
\]
这些标号由集合
\[
  I':=\{\alpha\beta,\beta\alpha\}\cup\{1,\ldots,k\}
\]
编号。注意 $Z(T')=Z(T)$。归纳假设断言，\BookRecordIndex{T0269}普适求值映射
\[
  \operatorname{ev}^{E'}:
  \mathcal M^*(\{A_\gamma\};\mathcal J^\ell)\times Z(T)
  \longrightarrow M^{E'}
\]
与 $\Delta^{E'}$ 横截。因此
\[
  \widetilde{\mathcal M}^*_{0,T'}(\{A_\alpha\};\mathcal J^\ell)
  =\bigl(\operatorname{ev}^{E'}\bigr)^{-1}(\Delta^{E'})
\]
是 $C^{\ell-1}$ \BookRecordIndex{T0053}巴拿赫流形。考虑\BookRecordIndex{T0283}求值映射
\[
  \operatorname{ev}_{\alpha\beta}:
  \widetilde{\mathcal M}^*_{0,T'}(\{A_\alpha\};\mathcal J^\ell)
  \longrightarrow M\times M,
\]
定义为
\[
  \operatorname{ev}_{\alpha\beta}(\mathbf u,\mathbf z,J)
  :=\bigl(u_\alpha(z_{\alpha\beta}),u_\beta(z_{\beta\alpha})\bigr).
\]
由定理~\ref{thm:6-3-1}，$M\times M$ 中每一点都是
$\operatorname{ev}_{\alpha\beta}$ 的\BookRecordIndex{T0409}正则值。因此，
$\operatorname{ev}_{\alpha\beta}$ 与对角线 $\Delta\subset M\times M$ 横截。由下面的
习题~\ref{ex:6-3-2}，可知
\[
  \operatorname{ev}^E
  =\operatorname{ev}^{E'}\times\operatorname{ev}_{\alpha\beta}
  \quad\text{与}\quad
  \Delta^E=\Delta^{E'}\times\Delta
\]
横截。这完成了归纳，也证明了命题~\ref{prop:6-2-8}。
\end{proof}

\begin{exercise}\label{ex:6-3-2}
设 $f_0:X\to Y_0$ 和 $f_1:X\to Y_1$ 是 $C^\ell$ \BookRecordIndex{T0053}巴拿赫流形之间的
$C^\ell$ 映射，而 $Z_0\subset Y_0$ 与 $Z_1\subset Y_1$ 是 $C^\ell$ 子流形。假设
$f_0$ 与 $Z_0$ 横截，并且 $f_1$ 在 $f_0^{-1}(Z_0)$ 上的限制与 $Z_1$ 横截。证明
$f_0\times f_1$ 与 $Z_0\times Z_1$ 横截。
\end{exercise}

考虑一棵有一个顶点和一个\BookRecordIndex{T0066}标记点的树。设
\[
  \operatorname{ev}:\mathcal M^*_{0,1}(A;\mathcal J)\longrightarrow M
\]
是相应的（光滑）\BookRecordIndex{T0269}普适求值映射。虽然这个\BookRecordIndex{T0283}求值映射总是浸没，但并不能由此\BookRecordIndex{T0326}推出存在某个\BookRecordIndex{T0088}殆复
结构 $J\in\mathcal J$，使\BookRecordIndex{T0283}求值映射
\[
  \operatorname{ev}_J:\mathcal M^*_{0,1}(A;J)\longrightarrow M
\]
是浸没。例如，\BookRecordIndex{T0246}模空间 $\mathcal M^*_{0,1}(A;J)$ 的维数可能就是太小。在下一个习题的
情形中还可以构造其他例子。该习题也说明了为何定理~\ref{thm:6-2-6} 必须限制到\BookRecordIndex{T0184}简单稳定
映射。它用到如下事实：若 $E$ 是
$\mathbb{CP}^2\#\overline{\mathbb{CP}}^{\,2}$ 中例外除子的类，则对通有的\BookRecordIndex{T0348}相容\BookRecordIndex{T0088}殆复结构
$J$，\BookRecordIndex{T0246}模空间
$\overline{\mathcal M}_{0,0}(E;J)$ 恰含一个元素（见例 7.1.15）。因此，由\BookRecordIndex{T0186}交点正性，
类 $2E$ 没有单纯代表。

\begin{exercise}\label{ex:6-3-3}
证明：对通有的 $J$，
\[
  \overline{\mathcal M}_{0,0}
  (\mathbb{CP}^2\#\overline{\mathbb{CP}}^{\,2},2E;J)
\]
是 $4$ 维流形，而它的虚维数为 $2$。当 $T$ 有两个顶点且没有\BookRecordIndex{T0066}标记点时，计算\BookRecordIndex{T0072}层
\[
  \mathcal M_{0,T}
  (\mathbb{CP}^2\#\overline{\mathbb{CP}}^{\,2},2E;J)
\]
的实际维数，并与虚维数比较。
\end{exercise}
